\section{Límite superior e inferior}

\begin{theorem}
  Sea ${(x_n)}_{n \in \N}$ una sucesión acotada. Entonces $\liminf x_n$ es el menor punto de acumulación de la sucesión y $\limsup x_n$ es el mayor.
\end{theorem}

\begin{proof}
  Sea $a = \liminf_{n \to \infty} x_n = \lim_{n \to \infty} a_n$, con $a_n = \inf X_n$.

  \emph{$a$ es punto de acumulación.} Sean $\e > 0$ y $n_0 \in \N$. Como $a_n \to a$, existe $n_1 > n_0$ con $a - \e < a_{n_1} \leq a$. Como $a_{n_1} = \inf X_{n_1}$, el número $a + \e > a_{n_1}$ no es cota inferior de $X_{n_1}$, luego existe $n \geq n_1$ con $x_n < a + \e$; además $x_n \geq a_{n_1} > a - \e$. Entonces $x_n \in (a - \e, a + \e)$ con $n > n_0$. Como $n_0$ es arbitrario, hay infinitos índices con $x_n \in (a - \e, a + \e)$, y por el teorema de la sección anterior, $a$ es punto de acumulación.

  \emph{Ningún $c < a$ es punto de acumulación.} Como $a_n \to a$ y $c < a$, existe $n_0$ con $c < a_{n_0} \leq a$. Sea $\e = a_{n_0} - c > 0$. Para $n \geq n_0$ se cumple $x_n \geq a_{n_0} = c + \e$, de modo que $x_n \notin (c - \e, c + \e)$ salvo quizás para $n < n_0$ (finitos índices). Luego $c$ no es punto de acumulación.

  La afirmación sobre $\limsup$ se prueba de manera análoga.
\end{proof}

\begin{theorem}[Bolzano--Weierstrass]
  Toda sucesión acotada de $\R$ tiene una subsucesión convergente.
\end{theorem}

\begin{proof}
  $\limsup x_n$ es punto de acumulación de ${(x_n)}$, es decir, el límite de alguna subsucesión.
\end{proof}

\begin{corollary}
  Una sucesión acotada de $\R$ es convergente si y sólo si $\limsup x_n = \liminf x_n$, es decir, si posee un único punto de acumulación.
\end{corollary}

\begin{proof}
  ($\Rightarrow$) Si ${(x_n)}$ converge a $L$, toda subsucesión converge a $L$. Como $\liminf x_n$ y $\limsup x_n$ son límites de subsucesiones, ambos son iguales a $L$.

  ($\Leftarrow$) Sea $a = \liminf x_n = \limsup x_n = \lim a_n = \lim b_n$. Dado $\e > 0$, existe $n_0$ tal que $a - \e < a_n$ y $b_n < a + \e$ para $n \geq n_0$. Como $x_n \in X_n$, se tiene $a_n \leq x_n \leq b_n$, y entonces $a - \e < x_n < a + \e$ para todo $n \geq n_0$. Es decir, $x_n \to a$.
\end{proof}

\section{Sucesiones de Cauchy}

Una sucesión ${(x_n)}_{n \in \N} \subset \R$ es \emph{de Cauchy} si para todo $\e > 0$ existe $n_0 \in \N$ tal que $|x_n - x_m| < \e$ para todos $n, m > n_0$, tal como se hizo en $\Q$. Valen, con las mismas demostraciones que en $\Q$:
\begin{enumerate}[label=(\arabic*)]
  \item Toda sucesión convergente es de Cauchy.
  \item Toda sucesión de Cauchy es acotada.
\end{enumerate}

\begin{lemma}
  Si una sucesión de Cauchy tiene una subsucesión que converge a $a \in \R$, entonces $\lim_{n \to \infty} x_n = a$.
\end{lemma}

\begin{proof}
  Dado $\e > 0$, existe $n_0$ tal que $|x_n - x_m| < \frac{\e}{2}$ para $n, m > n_0$. Como $a$ es el límite de una subsucesión, existe un índice $n_1 > n_0$ de la subsucesión con $|x_{n_1} - a| < \frac{\e}{2}$. Entonces, para $n > n_0$,
  \[
    |x_n - a| \leq |x_n - x_{n_1}| + |x_{n_1} - a| < \frac{\e}{2} + \frac{\e}{2} = \e. \qedhere
  \]
\end{proof}

\begin{theorem}
  Toda sucesión de Cauchy en $\R$ es convergente.
\end{theorem}

\begin{proof}
  Por ser de Cauchy es acotada; por el teorema de Bolzano--Weierstrass, alguna subsucesión converge, y por el lema anterior toda la sucesión converge al mismo límite.
\end{proof}

\section{Límites infinitos}

\begin{definition}
  Decimos que $\lim_{n \to \infty} x_n = +\infty$ si para todo $A > 0$ existe $n_0 \in \N$ tal que $x_n > A$ para todo $n > n_0$. Decimos que $\lim_{n \to \infty} x_n = -\infty$ si para todo $A > 0$ existe $n_0 \in \N$ tal que $x_n < -A$ para todo $n > n_0$.
\end{definition}

\begin{prop}
  Si $\lim_{n \to \infty} x_n = +\infty$ e ${(y_n)}$ es acotada inferiormente, entonces $\lim_{n \to \infty} (x_n + y_n) = +\infty$.
\end{prop}

\begin{proof}
  Sea $c \in \R$ con $c < y_n$ para todo $n$. Dado $A > 0$, existe $n_0$ tal que $x_n > A - c$ para $n > n_0$. Entonces $x_n + y_n > (A - c) + c = A$ para $n > n_0$.
\end{proof}

\begin{prop}
  Si $\lim_{n \to \infty} x_n = +\infty$ y existe $c > 0$ tal que $y_n > c$ para todo $n$, entonces $\lim_{n \to \infty} x_n y_n = +\infty$.
\end{prop}

\begin{proof}
  Dado $A > 0$, existe $n_0$ tal que $x_n > \frac{A}{c}$ para $n > n_0$. Entonces $x_n y_n > \frac{A}{c} \cdot c = A$ para $n > n_0$.
\end{proof}

\begin{prop}
  Si $x_n > 0$ para todo $n$, entonces $\lim_{n \to \infty} x_n = 0$ si y sólo si $\lim_{n \to \infty} \frac{1}{x_n} = +\infty$.
\end{prop}

\begin{proof}
  ($\Rightarrow$) Dado $A > 0$, existe $n_0$ tal que $0 < x_n < \frac{1}{A}$ para $n > n_0$, es decir, $\frac{1}{x_n} > A$.

  ($\Leftarrow$) Dado $\e > 0$, existe $n_0$ tal que $\frac{1}{x_n} > \frac{1}{\e}$ para $n > n_0$, es decir, $0 < x_n < \e$.
\end{proof}

\begin{prop}
  Sean ${(x_n)}$ e ${(y_n)}$ sucesiones de términos positivos.
  \begin{enumerate}[label=(\arabic*)]
    \item Si existe $c > 0$ con $x_n > c$ para todo $n$ y $\lim_{n \to \infty} y_n = 0$, entonces $\lim_{n \to \infty} \frac{x_n}{y_n} = +\infty$.
    \item Si ${(x_n)}$ es acotada y $\lim_{n \to \infty} y_n = +\infty$, entonces $\lim_{n \to \infty} \frac{x_n}{y_n} = 0$.
  \end{enumerate}
\end{prop}

\begin{proof}
  (1) Dado $A > 0$, existe $n_0$ tal que $0 < y_n < \frac{c}{A}$ para $n > n_0$, de modo que $\frac{x_n}{y_n} > c \cdot \frac{A}{c} = A$.

  (2) Sea $k > 0$ una cota superior de ${(x_n)}$. Dado $\e > 0$, existe $n_0$ tal que $y_n > \frac{k}{\e}$ para $n > n_0$, de modo que $0 < \frac{x_n}{y_n} < k \cdot \frac{\e}{k} = \e$.
\end{proof}

\section{Series numéricas}

Dada una sucesión ${(a_n)}_{n \in \N} \subset \R$ formamos la sucesión ${(S_n)}_{n \in \N}$ de \emph{sumas parciales} dada por $S_n = a_1 + \cdots + a_n$. La \emph{serie} $\sum_{n \geq 1} a_n$ \emph{converge} si ${(S_n)}$ converge, y su suma es $\lim_{n \to \infty} S_n$.

\begin{theorem}
  Si $\sum_{n \geq 1} a_n$ converge, entonces $\lim_{n \to \infty} a_n = 0$.
\end{theorem}

\begin{proof}
  Sea $s = \lim_{n \to \infty} S_n$. Como $S_{n-1} \to s$ también,
  \[
    a_n = S_n - S_{n-1} \longrightarrow s - s = 0. \qedhere
  \]
\end{proof}

\begin{eg}
  El recíproco del teorema anterior es falso: $\sum_{n \geq 1} \frac{1}{n}$ diverge.
\end{eg}

\begin{proof}
  Agrupando los términos en bloques de $2^{k-1}$ sumandos, cada uno de los cuales es al menos $2^{k-1} \cdot \frac{1}{2^k} = \frac{1}{2}$,
  \begin{align*}
    S_{2^n} & = 1 + \frac{1}{2} + \left( \frac{1}{3} + \frac{1}{4} \right) + \cdots + \left( \frac{1}{2^{n-1} + 1} + \cdots + \frac{1}{2^n} \right) \\
            & \geq 1 + \frac{1}{2} + \frac{2}{4} + \cdots + \frac{2^{n-1}}{2^n} = 1 + \frac{n}{2}.
  \end{align*}
  Luego $S_{2^n} \to +\infty$ y la serie diverge.
\end{proof}

\begin{theorem}
  Si $a_n \geq 0$ para todo $n$, entonces $\sum_{n \geq 1} a_n$ converge si y sólo si la sucesión ${(S_n)}$ es acotada.
\end{theorem}

\begin{proof}
  Por ser $a_n \geq 0$, ${(S_n)}$ es creciente, y por lo tanto converge si y sólo si es acotada.
\end{proof}

\begin{theorem}[Criterio de Cauchy]
  $\sum a_n$ converge si y sólo si para todo $\e > 0$ existe $n_0 \in \N$ tal que
  \[
    |a_{n+1} + \cdots + a_{n+p}| < \e \quad \text{para todo } n > n_0 \text{ y todo } p \in \N.
  \]
\end{theorem}

\begin{proof}
  Como $|S_{n+p} - S_n| = |a_{n+1} + \cdots + a_{n+p}|$, la condición dice exactamente que ${(S_n)}$ es de Cauchy, lo cual equivale a que sea convergente.
\end{proof}

\begin{theorem}[Criterio de comparación]
  Sean $a_n, b_n \geq 0$ para todo $n$ y supongamos que existen $c > 0$ y $n_0 \in \N$ con $a_n \leq c\, b_n$ para todo $n > n_0$.
  \begin{enumerate}[label=(\alph*)]
    \item Si $\sum b_n$ converge, entonces $\sum a_n$ converge.
    \item Si $\sum a_n$ diverge, entonces $\sum b_n$ diverge.
  \end{enumerate}
\end{theorem}

\begin{proof}
  Notemos que la convergencia de una serie sólo depende de su cola: $\sum a_n$ converge si y sólo si $\sum_{n > n_1} a_n$ converge para algún (todo) $n_1$.

  (a) Si $\sum b_n$ converge, sus sumas parciales $T_n$ están acotadas, digamos por $K$. Para $n > n_0$ las sumas parciales de $\sum a_n$ cumplen
  \[
    S_n = S_{n_0} + \sum_{m = n_0 + 1}^{n} a_m \leq S_{n_0} + c \sum_{m = n_0 + 1}^{n} b_m \leq S_{n_0} + c K.
  \]
  Luego ${(S_n)}$ es acotada y $\sum a_n$ converge.

  (b) Es la contrarrecíproca de (a).
\end{proof}

\begin{definition}
  La serie $\sum a_n$ es \emph{absolutamente convergente} si $\sum |a_n|$ converge. Si $\sum a_n$ converge pero $\sum |a_n|$ diverge, se dice que es \emph{condicionalmente convergente}.
\end{definition}

\begin{theorem}
  Toda serie absolutamente convergente es convergente.
\end{theorem}

\begin{proof}
  Dado $\e > 0$, como $\sum |a_n|$ converge, por el criterio de Cauchy existe $n_0$ tal que $|a_{n+1}| + \cdots + |a_{n+p}| < \e$ para $n > n_0$ y $p \in \N$. Entonces
  \[
    |a_{n+1} + \cdots + a_{n+p}| \leq |a_{n+1}| + \cdots + |a_{n+p}| < \e
  \]
  para $n > n_0$ y $p \in \N$, y por el criterio de Cauchy, $\sum a_n$ converge.
\end{proof}

\begin{theorem}[Criterio de condensación de Cauchy]
  Sea ${(a_n)}$ una sucesión con $a_1 \geq a_2 \geq \cdots \geq 0$. Entonces $\sum_{n \geq 1} a_n$ converge si y sólo si $\sum_{k \geq 0} 2^k a_{2^k}$ converge.
\end{theorem}

\begin{proof}
  Sean $s_n = a_1 + \cdots + a_n$ y $t_k = a_1 + 2 a_2 + \cdots + 2^k a_{2^k}$.

  Si $n < 2^k$, agrupando en bloques y usando que la sucesión es decreciente,
  \begin{align*}
    s_n & \leq a_1 + (a_2 + a_3) + \cdots + (a_{2^{k-1}} + \cdots + a_{2^k - 1}) \\
        & \leq a_1 + 2 a_2 + \cdots + 2^{k-1} a_{2^{k-1}} \leq t_k.
  \end{align*}
  Por otro lado, si $n \geq 2^k$,
  \begin{align*}
    s_n & \geq a_1 + a_2 + (a_3 + a_4) + \cdots + (a_{2^{k-1} + 1} + \cdots + a_{2^k}) \\
        & \geq a_1 + a_2 + 2 a_4 + \cdots + 2^{k-1} a_{2^k} \geq \frac{1}{2} t_k.
  \end{align*}
  Luego las sucesiones ${(s_n)}$ y ${(t_k)}$ (crecientes) son ambas acotadas o ambas no acotadas, y el resultado se sigue del teorema sobre series de términos no negativos.
\end{proof}
